% Fault101 Fault 2_2 Explorer - Step-by-Step Student Guide
% Generated by guide_tools/latex/latexkit.py from the story script; this file is plain LaTeX - edit it directly.
% Build: pdflatex <this file>   (run it twice so the contents page fills in).  Needs scaguide.sty and figures/ next to it.
\documentclass[11pt,letterpaper]{article}
\usepackage{scaguide}
\guidesetup{
  headerleft={ChipWhisperer Fault101  -  Fault 2\_2 Explorer Guide},
  headerright={Password Bypass Explorer},
  footerleft={Companion to: Fault 2\_2 - Voltage Glitching to Bypass Password Explorer.ipynb},
  title={Fault101 Fault 2\_2 Explorer - Step-by-Step Student Guide},
  subject={Step-by-step guide to the Fault 2\_2 Voltage Glitching password-bypass Explorer: why Fault 2\_1's settings do not carry over, and how to search for ones that do}}

\begin{document}

\guidecover{FAULT101  -  FAULT 2\_2  (EXPLORER)}{Password Bypass\\ Explorer}{Step-by-step student guide: searching for the fault that makes a wrong password pass}

Fault 2\_2 asks the same question as Fault 2\_1, against different firmware: can a brief voltage glitch make a check
come out wrong? Here the check is a password comparison instead of a counting loop. This guide's biggest lesson
arrives early: \textbf{the settings that worked perfectly for Fault 2\_1 do almost nothing useful here} - they just
crash the target. Finding a setting that actually flips the password check's answer means widening the search, and
even then, what you find will not sit still: these settings \textbf{drift} over the course of a session. \textbf{This
lab needs real hardware; there is no simulated mode.}

\vspace{10pt}

\begin{guidecallout}{tip}{HOW TO USE THIS GUIDE}
This guide goes \textbf{with} the notebook - it does not replace it. Keep the notebook open next to this PDF.
\end{guidecallout}

\guidehthree{\textbf{Which files are which}}

{\small\arrayrulecolor{guidegrid}
\rowcolors{2}{white}{guidezebra}
\begin{longtable}{|R{0.4348}|R{0.4493}|R{0.1159}|}
\hline
\rowcolor{guidenavy}\hdrcell{File (folder: jupyter/courses/fault101/)} & \hdrcell{What it is} & \hdrcell{Open it?} \\ \hline
\endhead
\codesmall{Fault 2\_\bk{}2 - Voltage Glitching to Bypass Password Explorer.\bk{}ipynb} & The Explorer. Everything in this guide happens here. & \textbf{Yes} \\ \hline
\codesmall{Fault 2\_\bk{}2 - Voltage Glitching to Bypass Password.\bk{}ipynb} & The original (non-Explorer) lab. & Optional \\ \hline
\codesmall{SOLN\_\bk{}Fault 2\_\bk{}2 - .\bk{}.\bk{}.\bk{}ipynb} & Answer key for the original lab. & If stuck \\ \hline
\codesmall{firmware/\bk{}mcu/\bk{}simpleserial-glitch/\bk{}simpleserial-glitch.\bk{}c} & (a different folder: \codesmall{chipwhisperer/\bk{}firmware/\bk{}.\bk{}.\bk{}.\bk{}}) The \codesmall{password()} function this lab targets. \textbf{Step 8 edits the \codesmall{passwd[]} line of it.} & Step 8 \\ \hline
\end{longtable}}

\vspace{4pt}

\guidehthree{\textbf{The four colored boxes you will see everywhere}}

\guidelegend

\vspace{4pt}

\guidesmall{Guide version 1.0 - November 2026. Every number in this guide was measured on one real kit:
ChipWhisperer-Husky, a CW313 board and a CW312T-SAM4S target (Windows 11, Python 3.12, ChipWhisperer 6.0.0, firmware
1.5.0) - the same kit as the Fault 2\_1 guide. \textbf{Read the drift warning after Experiment 5 before you trust any
single number in this guide}: several of the exact figures below changed - in one case from 15\% success to 0\% - when
we re-measured them only minutes later in the same session.}

\clearpage

\tableofcontents

\clearpage

\guidehone{The big picture}

\textbf{What is different from Fault 2\_1?} The target firmware now runs \code{password()}: it compares a 5-byte guess
against a built-in constant (\code{\textquotedbl{}touch\textquotedbl{}} - a fixed string baked into this open-source
demo firmware, not a secret worth hiding) and reports pass or fail. Unlike the SCA101 password lab, this check never
stops early: it always runs all five comparisons, so there is no timing information to measure - the \textbf{only} way
to make a wrong guess pass is to disturb the chip while it runs, corrupting the comparison's result directly.

\textbf{Why doesn't Fault 2\_1's setting just work again?} A glitch setting disturbs *whatever instruction is
executing at that instant*. The counting loop and the password check are different code, running at different times,
so the instant that mattered for one means nothing for the other. This Explorer makes you discover that the hard way
before showing you how to search past it.

\guidehthree{Your path through the notebook:}

\begin{center}
\begin{tikzpicture}
\node[rounded corners=4pt,fill=guidenavy,text=white,align=center,minimum height=1.15cm,minimum width=2.46cm,inner sep=3pt] (n0) at (1.23,0) {\textbf{1  Reuse}\\{\footnotesize Fault 2\_1's setting}};
\node[rounded corners=4pt,fill=guidenavy,text=white,align=center,minimum height=1.15cm,minimum width=2.46cm,inner sep=3pt] (n1) at (4.24,0) {\textbf{2 - 3  Map}\\{\footnotesize the nearby area}};
\node[rounded corners=4pt,fill=guidenavy,text=white,align=center,minimum height=1.15cm,minimum width=2.46cm,inner sep=3pt] (n2) at (7.25,0) {\textbf{4  Search}\\{\footnotesize a wider box}};
\node[rounded corners=4pt,fill=guidenavy,text=white,align=center,minimum height=1.15cm,minimum width=2.46cm,inner sep=3pt] (n3) at (10.25,0) {\textbf{5 - 6  Zoom}\\{\footnotesize in and compare}};
\node[rounded corners=4pt,fill=guidenavy,text=white,align=center,minimum height=1.15cm,minimum width=2.46cm,inner sep=3pt] (n4) at (13.26,0) {\textbf{7  Check}\\{\footnotesize the guess itself}};
\node[rounded corners=4pt,fill=guideteal,text=white,align=center,minimum height=1.15cm,minimum width=2.46cm,inner sep=3pt] (n5) at (16.27,0) {\textbf{8  Choose}\\{\footnotesize your own password}};
\draw[-{Stealth},guidegrey,thick] (n0.east) -- (n1.west);
\draw[-{Stealth},guidegrey,thick] (n1.east) -- (n2.west);
\draw[-{Stealth},guidegrey,thick] (n2.east) -- (n3.west);
\draw[-{Stealth},guidegrey,thick] (n3.east) -- (n4.west);
\draw[-{Stealth},guidegrey,thick] (n4.east) -- (n5.west);
\end{tikzpicture}
\end{center}

\vspace{6pt}

{\small\arrayrulecolor{guidegrid}
\rowcolors{2}{white}{guidezebra}
\begin{longtable}{|R{0.1159}|R{0.3913}|R{0.4928}|}
\hline
\rowcolor{guidenavy}\hdrcell{Experiment} & \hdrcell{What you do} & \hdrcell{What you learn} \\ \hline
\endhead
\textbf{1} & Fire Fault 2\_1's exact winning setting at this firmware. & It does not bypass anything - it just crashes the target. \\ \hline
\textbf{2} & Scan width to map this firmware's own crash edge. & A real crash boundary exists here too, at a similar width - but no success hides next to it. \\ \hline
\textbf{3} & Scan \codesmall{ext\_offset} widely at that edge. & A periodic crash pattern, but still no success anywhere nearby. \\ \hline
\textbf{4} & Run a randomized search across a much bigger box. & Success lives in a completely different corner of the search space - and how often you find one varies a lot between runs. \\ \hline
\textbf{5} & Zoom in on a success and map its neighbourhood. & A usable plateau in width/offset - and a hard lesson: these settings drift over time. \\ \hline
\textbf{6} & Compare \codesmall{hp} alone against \codesmall{hp + lp} together. & Stronger is not always better - and even that comparison needs re-checking, because of the drift. \\ \hline
\textbf{7} & Check whether the guess itself matters. & It does not: this is a control-flow fault, not a side channel. \\ \hline
\textbf{8} & Change the firmware's password to your own, rebuild, reflash, and glitch it. & The same setting that worked against ``touch'' keeps working, at roughly the same rate, against a password you chose yourself. \\ \hline
\end{longtable}}

\clearpage

\guidehone{Before you start}

Same kit and checklist as Fault 2\_1.

{\small\arrayrulecolor{guidegrid}
\rowcolors{2}{white}{guidezebra}
\begin{longtable}{|R{0.0652}|R{0.5217}|R{0.4130}|}
\hline
\rowcolor{guidenavy}\hdrcell{} & \hdrcell{You need} & \hdrcell{How to check} \\ \hline
\endhead
[  ] & Fault 2\_1 completed (same setup, same hardware checklist). & You can fire a glitch and see \codesmall{normal}/\codesmall{reset}/\codesmall{success} outcomes. \\ \hline
[  ] & \textbf{ChipWhisperer-Husky} plugged straight into a USB port (not through a hub). & Lab 0's hardware test ends with \codesmall{hardware OK (disconnected again)}. \\ \hline
[  ] & \textbf{CW312T-SAM4S} target seated on the \textbf{CW313} board, ribbon connected. & See the SCA101 Linux guide, Part 4.2. \\ \hline
[  ] & The Husky \textbf{\codesmall{Crowbar}} SMA cable connected to the CW313's second SMA socket. & Without it, every glitch here will show as \codesmall{normal} too. \\ \hline
[  ] & No other notebook holding the Husky open. & \textbf{File \textgreater{} Shut Down Notebook} in any other open notebook. \\ \hline
\end{longtable}}

\vspace{4pt}

\begin{guidecallout}{warn}{WATCH OUT}
\textbf{This lab can take longer than Fault 2\_1.} Experiment 4's randomized search ran for up to 4 minutes on our kit
before finding a working setting, and sometimes found nothing at all in that time. That is expected - plan for it, and
do not assume something is broken just because a search does not finish quickly.
\end{guidecallout}

\guidehone{Experiment 1 - try Fault 2\_1's winning setting here}

\textbf{Goal:} find out whether the setting that reliably glitched Fault 2\_1's counting loop (\code{WIDTH=1770,
OFFSET=2670, EXT\_\bk{}OFFSET=1}) does anything useful against this firmware's password check.

\textbf{Do this:} run the sanity-check cell first (it confirms the correct password \code{touch} passes and wrong
guesses fail, with no glitching). Then run the Experiment 1 cell with a \textbf{wrong} guess at Fault 2\_1's setting.

\begin{guidecallout}{see}{YOU SHOULD SEE}
\textbf{Almost every shot is \code{reset}.} The setting that calmly disturbed a 2500-count loop is far too strong for
whatever the chip is doing at that same instant while checking a password - it simply crashes the target, every time,
regardless of which wrong guess you send. A glitch setting belongs to *one piece of code at one moment*, not to the
chip in general.
\end{guidecallout}

\guidehone{Experiment 2 - map this firmware's own crash edge}

\textbf{Goal:} Experiment 1 sat deep inside a crash region. Find where that region's edge actually is for
\textbf{this} firmware, the same way Fault 2\_1 Experiment 2 did for the counting loop.

\begin{enumerate}[label=\arabic*.,leftmargin=1.8em,itemsep=2pt,topsep=2pt,parsep=0pt,start=1]
\item Run the cell as it is: it scans width at a fixed offset and ext\_offset, with a wrong guess.
\item Change \code{EXT\_OFFSET} to 2, then to 0. Does the crash region change?
\item Zoom in with \code{WIDTH\_STEP = 2} around the edge you find.
\end{enumerate}

\begin{guidecallout}{see}{YOU SHOULD SEE}
At \code{ext\_offset=1} (offset 2670): width 1750 and below is \code{normal}; 1758-1764 is a mix of \code{normal} and
\code{reset}; 1766 and above is \code{reset} every time. \code{ext\_offset=2} gives a very similar edge a couple of
steps earlier. \code{ext\_offset=0}, and every value from 3 up to about 5, never disturbs the target at all across
this whole width range - only cycles 1 and 2 are sensitive here, exactly like Fault 2\_1's counting loop. \textbf{No
width in this scan ever produced a \code{success}.} Crashing this check is easy; making it lie from here is not.
\end{guidecallout}

\guidehone{Experiment 3 - scan a wide range of moments}

\textbf{Goal:} this password check's trigger window stays high much longer than one or two clock cycles - it covers
five comparisons plus a serial write-back. Scan \code{ext\_offset} widely to see the rest of that window.

\textbf{Do this:} run the cell as it is - a width just below Experiment 2's edge, scanned across \code{ext\_offset} 0
to 60.

\begin{guidecallout}{see}{YOU SHOULD SEE}
\code{reset} appears in short bursts roughly every 20 cycles (around 1-2, 20-25, 33, 40-42, 53 and 59-60), with
\code{normal} everywhere else - probably the serial/UART handling inside the trigger window, not the comparison
itself. \textbf{Still no \code{success} anywhere in this range.} We have now swept 60 different moments at a width
right on an edge, and found nothing: whatever works here, it is not a small tweak away from where Fault 2\_1's pattern
would suggest looking.
\end{guidecallout}

\guidehone{Experiment 4 - widen the net: a randomized search}

\textbf{Goal:} three experiments in the `obvious' area found nothing but crashes. Rather than inching the ranges wider
by hand, let the computer try \textbf{random} combinations across a much bigger box and watch for the first success.

\textbf{Do this:} run the cell as it is. \textbf{This can take a few minutes} and may find nothing within the time
budget - both are normal outcomes here. For a sense of scale: our test runs resolved anywhere from under a minute to
the full 4-minute default budget. \textbf{If you are running this live with a class and need a hard time limit, lower
\code{TIME\_BUDGET\_S}} - just expect the `no bypass found' outcome to come up more often, which is still a real
result worth discussing, not a failure to fix.

\begin{guidecallout}{see}{YOU SHOULD SEE}
This is the one experiment where how much luck you need varies a lot between runs. In one run, a search like this
tried \textbf{7248 settings over about 230 seconds and found only 2 successes} - roughly 1 in 3600 attempts. Run again
later in the same session, it found one in \textbf{1149 attempts, about 46 seconds} - a very different hit rate (see
the drift warning after Experiment 5: the *size* of the good region itself seems to change over time, not just its
exact location). What stayed consistent across every run: successes shared a \textbf{much bigger width (close to
1970-1990) and a much bigger offset (close to 3000-3100)} than anywhere in Experiments 1 to 3 - a completely different
corner of the search space from the one that merely crashes the target.
\end{guidecallout}

\begin{guidecallout}{tip}{GOOD TO KNOW}
\textbf{If your run finds nothing in the time budget, that is a real result, not a failure.} It means this particular
fault, from this glitch circuit, did not defeat this check within the search you tried - not every fault-and-target
combination gives in, and recording that honestly is part of a real security evaluation.
\end{guidecallout}

\guidehone{Experiment 5 - zoom in on what you found}

\textbf{Goal:} turn one lucky random hit into a setting you understand, and see how wide the `good' neighbourhood
actually is.

\begin{enumerate}[label=\arabic*.,leftmargin=1.8em,itemsep=2pt,topsep=2pt,parsep=0pt,start=1]
\item Put the \code{(width, offset, ext\_\bk{}offset)} Experiment 4 found into \code{CENTER} (or start from ours: \code{(1979, 3071, 121)}).
\item Run the cell: it tries a small grid of neighbours around \code{CENTER}, all at the same \code{ext\_offset}.
\item Change \code{EXT\_OFFSET\_TRY} to a few cycles either side of \code{CENTER}`s. Is the neighbourhood just as wide there?
\end{enumerate}

\begin{guidecallout}{see}{YOU SHOULD SEE}
When we first mapped it, most of a 20-by-40-step block around \code{(1979, 3071)} gave at least one success out of 8
shots, with individual points ranging from about 10\% to 40\% - a genuinely usable plateau, much wider than the thin
diagonal line Fault 2\_1 found. \code{ext\_offset} was pickier: only a few specific values (121 and 130 among the ones
we tried, and 89 at a different width/offset) gave any success at all there.
\end{guidecallout}

\begin{guidecallout}{warn}{WATCH OUT}
\textbf{These settings drift.} When we came back to the *exact same* \code{(1979, 3071, 121)} later in the same
session, it gave \textbf{zero} successes in 60 shots - not one. A nearby point our random search found independently,
\code{(1973, 3083, 141)}, was still giving occasional successes at that point. Nothing about the wiring or the code
had changed; the glitch circuit's precise sweet spot moves over time, likely with the board's temperature and small
clock drift, by more than you might expect. \textbf{Treat every specific number in this guide as `this was possible,'
not as a guarantee you can reproduce on demand.} Always re-measure with fresh shots before trusting a setting, and if
one stops working, try its near neighbours before assuming it never worked at all.
\end{guidecallout}

\guidehone{Experiment 6 - does a stronger glitch work better here too?}

\textbf{Goal:} Fault 2\_1 found that \code{hp + lp} together crashed the target at a smaller width than \code{hp}
alone - a stronger glitch. Does using both transistors also help bypass \textbf{this} check?

\textbf{Do this:} run the cell at your best setting from Experiment 5, once with \code{hp} only and once with \code{hp
+ lp}.

\begin{guidecallout}{see}{YOU SHOULD SEE}
In one back-to-back comparison at a setting that was working at the time: \code{hp} alone gave about 15\% success (6
of 40), \code{hp + lp} together gave \textbf{0\%} - every shot was a \code{reset}. The stronger glitch did not help;
it pushed the disturbance past the useful instant straight into a crash - the opposite of Fault 2\_1's counting loop.
\textbf{The caveat:} because these settings drift, treat a single comparison like this as one data point, not a law -
re-run both with fresh shots at your own current best setting before you trust which one wins. What should not change
is the general lesson: which transistor combination helps depends on the code being glitched, not just on the physical
circuit, so it is always worth comparing both instead of assuming `stronger is always better'.
\end{guidecallout}

\guidehone{Experiment 7 - does the guess itself matter?}

\textbf{Goal:} measure your best setting's reliability, and check something important: does the *content* of the wrong
guess change anything? In the SCA101 power-analysis lab, the guess mattered enormously - that attack is a side
channel. If this fault works (or doesn't) regardless of what you send, that tells you it is a pure control-flow
corruption, not a leak of any kind.

\textbf{Do this:} run the cell as it is - it fires the same number of glitches at several different wrong guesses,
back-to-back at the same setting, and add a guess of your own.

\begin{guidecallout}{see}{YOU SHOULD SEE}
Run all the guesses back-to-back (so drift has the least chance to matter between them) and compare their success
counts \textbf{to each other}, not to a fixed target percentage. On our kit they came out close to each other,
whatever the absolute level happened to be at that moment - no guess stood out as easier or harder, including a guess
that was only one character away from the real password. That confirms the lesson from the firmware listing: the fault
corrupts the comparison's *outcome*, not a part of it that depends on which byte was wrong, so there is nothing to
gain from a `smarter' wrong guess - unlike a real side-channel attack, where the guess is everything.
\end{guidecallout}

\guidehone{Experiment 8 - choose your own password (hardware)}

Every experiment so far has glitched the same built-in password, \code{\textquotedbl{}touch\textquotedbl{}} - a line
in the firmware's source code, not a secret anyone is hiding. This step has you change it yourself, rebuild, reflash,
and glitch \textbf{your own} password, so you see firsthand - not just told - that the fault does not care what the
password is.

\guidehthree{\textbf{Where the password lives}}

{\small\arrayrulecolor{guidegrid}
\rowcolors{1}{guidezebra}{white}
\begin{longtable}{|R{0.1594}|R{0.8406}|}
\hline
\textbf{Folder} & \codesmall{firmware/\bk{}mcu/\bk{}simpleserial-glitch/\bk{}} - from the notebook, \codesmall{.\bk{}.\bk{}/\bk{}.\bk{}.\bk{}/\bk{}.\bk{}.\bk{}/\bk{}firmware/\bk{}mcu/\bk{}simpleserial-glitch/\bk{}} \\ \hline
\textbf{File} & \codesmall{simpleserial-glitch.\bk{}c} - the only file you edit \\ \hline
\textbf{Line} & \codesmall{char passwd[] = \textquotedbl{}touch\textquotedbl{};} inside \codesmall{password()} (search for \codesmall{passwd[]}) \\ \hline
\end{longtable}}

\vspace{4pt}

\guidehthree{\textbf{What counts as a good password here}}

{\small\arrayrulecolor{guidegrid}
\rowcolors{2}{white}{guidezebra}
\begin{longtable}{|R{0.3768}|R{0.6232}|}
\hline
\rowcolor{guidenavy}\hdrcell{Rule} & \hdrcell{Why} \\ \hline
\endhead
\textbf{Exactly 5 characters} & \codesmall{password()}`s loop always compares 5 bytes (\codesmall{for(cnt = 0; cnt \textless{} 5; cnt++)}). A shorter string would read past its own end; a longer one would just have its extra characters ignored. The notebook's \codesmall{use\_password()} enforces this. \\ \hline
Plain printable characters, no \codesmall{\textquotedbl{}} or \codesmall{\textbackslash{}} & Anything else either breaks the C string on that line or cannot be typed as a guess. \\ \hline
\end{longtable}}

\vspace{4pt}

\begin{guidecallout}{do}{DO THIS}
Run the helper cell once, set \code{MY\_PASSWORD} to your own 5-character string, and run
\code{use\_\bk{}password(MY\_\bk{}PASSWORD)}. It edits the file, rebuilds, reflashes, and checks the new password with
the chip for you.
\end{guidecallout}

\begin{guidecallout}{see}{YOU SHOULD SEE}
\code{passok=1} for your new password, \code{passok=0} for a wrong guess - and if you try the old
\code{\textquotedbl{}touch\textquotedbl{}}, that comes back \code{passok=0} too: the chip only knows the password you
just gave it.
\end{guidecallout}

\textbf{Now glitch it.} Put one of your own best settings from Experiments 5-7 into the next cell (or start from ours)
and fire a wrong guess at it, exactly as in the earlier experiments - no change to the attack itself, only the
firmware underneath it.

\begin{guidecallout}{see}{YOU SHOULD SEE}
\textbf{The same setting that worked against \code{\textquotedbl{}touch\textquotedbl{}} kept working, at roughly the
same rate, after we changed the password} - we measured it three separate times: 15\% (6 of 40) against
\code{\textquotedbl{}touch\textquotedbl{}}, then 17\% (5 of 30) and separately 10\% (3 of 30) against two different
runs with a freshly-chosen 5-character password. None of those three numbers match exactly, but they all sit in the
same rough range - exactly what you would expect given how much these settings already drift between runs of the
*same* password (Experiment 2). That makes sense once you remember why: changing
\code{\textquotedbl{}touch\textquotedbl{}} to another 5-character string does not change the shape of the compiled
comparison loop, only the constant it compares against - the same instant in the program is still the vulnerable one.
If your own run doesn't reproduce this at all, drift happens here too - use Experiment 4's search again; it works on
whatever password is currently flashed, unchanged.
\end{guidecallout}

\begin{guidecallout}{warn}{WATCH OUT}
\textbf{Put the original password back when you are done.} The last cell in this experiment has a
\code{RESTORE\_ORIGINAL} flag (default \code{False}, so nothing happens by accident) - set it to \code{True} and run
it to restore \code{\textquotedbl{}touch\textquotedbl{}} and reflash, so the next person starts from the same place
you did. The helper functions also take care to preserve the firmware file's exact line endings when they edit it, so
a careful restore leaves \code{git status} clean - check it yourself with \code{git status} in the \code{firmware/}
folder if you want to be sure.
\end{guidecallout}

\guidehone{Challenges (optional)}

\begin{enumerate}[label=\arabic*.,leftmargin=1.8em,itemsep=2pt,topsep=2pt,parsep=0pt,start=1]
\item \textbf{Characterize the second spot.} Map the neighbourhood around \code{(1985, 2999, 89)} the way Experiment 5 mapped \code{(1979, 3071, 121)}. Is it as wide?
\item \textbf{A tighter automatic search.} Now that you know roughly where the good region is, write a \code{GlitchController} search (like Fault 2\_1 Experiment 4) with \code{width} and \code{offset} ranges centred on it. How much faster is it than Experiment 4's blind random search?
\item \textbf{Does \code{repeat} help or hurt?} Try \code{REPEAT = 2} at your best setting. Does the success rate go up or down?
\item \textbf{A fourth spot.} Run Experiment 4 again (a new random seed) for several minutes. Does it always find settings with a large width and a large offset, or can a bypass exist somewhere else too?
\end{enumerate}

\guidehone{If something goes wrong}

{\small\arrayrulecolor{guidegrid}
\rowcolors{2}{white}{guidezebra}
\begin{longtable}{|R{0.2899}|R{0.3043}|R{0.4058}|}
\hline
\rowcolor{guidenavy}\hdrcell{What you see} & \hdrcell{Most likely cause} & \hdrcell{What to do} \\ \hline
\endhead
The sanity check does not print \codesmall{passok=1} for \codesmall{touch}. & Firmware did not flash, or a stale serial buffer. & Re-run the build/flash cells; re-run \codesmall{reboot\_flush()}. \\ \hline
Every shot is \codesmall{normal}, never \codesmall{reset} or \codesmall{success}, even at a huge width. & The Crowbar SMA cable is not connected. & Fit the cable (Fault 2\_1's guide), run the setup cells again. \\ \hline
Every shot is \codesmall{reset}. & You are inside a crash region for this width/offset. & Try a smaller width, or a different part of the search space (Experiment 4). \\ \hline
Experiment 4 finds nothing in the time budget. & A real result - see Experiment 4's `what you should see'. & Run it again (new random settings), or raise \codesmall{TIME\_BUDGET\_S}. \\ \hline
A setting that worked earlier stops working. & Drift (see the warning after Experiment 5). & Re-measure with fresh shots; try nearby settings too. \\ \hline
\codesmall{use\_password()} raises a \codesmall{ValueError}. & Your password is not exactly 5 characters, or uses \codesmall{\textquotedbl{}}, \codesmall{\textbackslash{}}, or a non-printable character. & Pick a different 5-character password and try again. \\ \hline
\codesmall{use\_password()} raises \codesmall{RuntimeError: firmware build failed}. & A typo elsewhere in \codesmall{simpleserial-glitch.\bk{}c}, or a stray edit from a previous session. & Read the build output printed above the error; if the file looks wrong, run \codesmall{restore\_\bk{}original\_\bk{}password()} and start Experiment 8 again. \\ \hline
\codesmall{git status} shows \codesmall{simpleserial-glitch.\bk{}c} changed after you restored it. & You restored the password but left \codesmall{RESTORE\_\bk{}ORIGINAL = False}, or edited the file by hand outside the notebook. & Run Experiment 8's last cell with \codesmall{RESTORE\_\bk{}ORIGINAL = True}, or \codesmall{git checkout -- firmware/\bk{}mcu/\bk{}simpleserial-glitch/\bk{}simpleserial-glitch.\bk{}c} from a terminal. \\ \hline
\codesmall{OSError: Could not find ChipWhisperer}. & Husky unplugged, or another notebook still holds it. & Plug it in; shut down every other notebook; restart this kernel. \\ \hline
\codesmall{USBErrorPipe} / \codesmall{Unknown hwInfoVer}. & The Husky's USB link or FPGA stopped answering. & Unplug and replug the Husky directly into the computer; restart the kernel. \\ \hline
A cell was interrupted (Stop button). & Target or glitch hardware may be left in an odd state. & Run \codesmall{recover()} in a new cell. \\ \hline
\end{longtable}}

\vspace{4pt}

\guidehone{Appendix - every variable you can edit}

{\small\arrayrulecolor{guidegrid}
\rowcolors{2}{white}{guidezebra}
\begin{longtable}{|R{0.2174}|R{0.1449}|R{0.6377}|}
\hline
\rowcolor{guidenavy}\hdrcell{Variable} & \hdrcell{Default} & \hdrcell{What it does} \\ \hline
\endhead
\codesmall{GUESS} & \codesmall{b\textquotedbl{}xxxxx\textquotedbl{}} & A 5-byte wrong guess sent to the password check. \\ \hline
\codesmall{WIDTH}, \codesmall{OFFSET}, \codesmall{EXT\_OFFSET} & varies & The glitch settings fired, in Husky phase-shift steps / clock cycles. \\ \hline
\codesmall{REPEAT} & 1 & Number of glitch pulses fired per shot. \\ \hline
\codesmall{SHOTS} & varies & How many glitches to fire for one measurement. \\ \hline
\codesmall{WIDTH\_\bk{}START/\bk{}STOP/\bk{}STEP} & varies & The width range scanned in Experiment 2. \\ \hline
\codesmall{EXT\_START/STOP/STEP} & 0 / 60 / 1 & The ext\_offset range scanned in Experiment 3. \\ \hline
\codesmall{WIDTH\_RANGE}, \codesmall{OFFSET\_RANGE}, \codesmall{EXT\_OFFSET\_RANGE} & wide & The box searched randomly in Experiment 4. \\ \hline
\codesmall{TIME\_BUDGET\_S} & 240 & How long Experiment 4 keeps trying before giving up. \\ \hline
\codesmall{CENTER} & a found success & The \codesmall{(width, offset, ext\_\bk{}offset)} Experiment 5 zooms in around. \\ \hline
\codesmall{WIDTH\_STEPS}, \codesmall{OFFSET\_STEPS} & lists & Offsets from \codesmall{CENTER} tried in Experiment 5's grid. \\ \hline
\codesmall{GUESSES} & list & The guesses compared in Experiment 7. \\ \hline
\codesmall{MY\_PASSWORD} & \codesmall{\textquotedbl{}f1zz9\textquotedbl{}} & The 5-character password Experiment 8 flashes onto the chip. \\ \hline
\codesmall{RESTORE\_ORIGINAL} & \codesmall{False} & Set to \codesmall{True} and run Experiment 8's last cell to put \codesmall{\textquotedbl{}touch\textquotedbl{}} back. \\ \hline
\end{longtable}}

\vspace{4pt}

\guidehone{Appendix - Glossary}

{\small\arrayrulecolor{guidegrid}
\rowcolors{2}{white}{guidezebra}
\begin{longtable}{|R{0.2464}|R{0.7536}|}
\hline
\rowcolor{guidenavy}\hdrcell{Term} & \hdrcell{Meaning} \\ \hline
\endhead
Voltage glitch / fault injection & See the Fault 2\_1 guide's glossary - the same concepts apply here. \\ \hline
Control-flow corruption & A fault that changes *what the program does* (skips a step, flips a result) rather than leaking information through timing or power, as a side-channel attack would. \\ \hline
Drift & The tendency of a glitch setting's success rate to change over time on the same kit, likely with temperature and small clock variations - not a measurement mistake. \\ \hline
Plateau & A region of settings (not just one point) that all give some success rate - as opposed to a single pinpoint setting. \\ \hline
\codesmall{passok} & The one-byte result the password check sends back: 1 for pass, 0 for fail. \\ \hline
\end{longtable}}

\vspace{4pt}

\guidehthree{Notice}

\guidesmall{This guide accompanies ChipWhisperer course material that references (commercial) training material on
ChipWhisperer.io. You can freely execute and use the lab per the open-source license (including in your own courses if
you distribute similarly), but you must maintain notice about the source: https://www.ChipWhisperer.io.}

\guidesmall{Copyright (C) NewAE Technology Inc., 2015-2020. ChipWhisperer is a trademark of NewAE Technology Inc.,
claimed in all jurisdictions, and registered in at least the United States of America, European Union, and Peoples
Republic of China. Tutorials derived from our open-source work must be released under the associated open-source
license, and notice of the source must be clearly displayed.}

\guidesmall{THE SOFTWARE IS PROVIDED ``AS IS'', WITHOUT WARRANTY OF ANY KIND, EXPRESS OR IMPLIED, INCLUDING BUT NOT
LIMITED TO THE WARRANTIES OF MERCHANTABILITY, FITNESS FOR A PARTICULAR PURPOSE AND NONINFRINGEMENT. IN NO EVENT SHALL
THE AUTHORS OR COPYRIGHT HOLDERS BE LIABLE FOR ANY CLAIM, DAMAGES OR OTHER LIABILITY, WHETHER IN AN ACTION OF
CONTRACT, TORT OR OTHERWISE, ARISING FROM, OUT OF OR IN CONNECTION WITH THE SOFTWARE OR THE USE OR OTHER DEALINGS IN
THE SOFTWARE.}

\end{document}
